% ECONOMICS_PROBLEM: {"id": "D63-86", "title": "Exact original maximin shares using at most one duplicated good", "jel_code": "D63", "source_id": "OP-0195"}
% !TeX program = xelatex
\documentclass[11pt,a4paper]{article}
\usepackage[margin=23mm]{geometry}
\usepackage{fontspec}
\setmainfont{Latin Modern Roman}
\usepackage{amsmath,amssymb,mathtools}
\usepackage{xcolor,enumitem,booktabs,longtable,array,xurl,hyperref}
\definecolor{muted}{HTML}{53636C}
\hypersetup{colorlinks=true,urlcolor=blue,linkcolor=blue}
\setlength{\parindent}{0pt}
\setlength{\parskip}{5pt}
\newcommand{\R}{\mathbb R}
\newcommand{\E}{\mathbb E}
\newcommand{\N}{\mathbb N}
\newcommand{\ind}{\mathbf 1}
\newcommand{\Var}{\operatorname{Var}}
\newcommand{\Cov}{\operatorname{Cov}}
\newcommand{\supp}{\operatorname{supp}}
\DeclareMathOperator*{\argmin}{arg\,min}
\DeclareMathOperator*{\argmax}{arg\,max}
\newcommand{\entrymeta}[3]{{\sffamily\small\color{muted}\textbf{\texttt{#1}}\quad|\quad #2\par #3\par}}
\newcommand{\Problemref}[1]{\texttt{#1}}
\newcommand{\JELref}[2]{\texttt{#2} (JEL \texttt{#1})}
\begin{document}
\section*{Exact original maximin shares using at most one duplicated good}
\textbf{Problem ID:} \texttt{D63-86}\quad \textbf{JEL:} \texttt{D63}\par
% BEGIN ECONOMICS STATEMENT
\label{problem:OP-0195}
% GAME_THEORY_PROBLEM: {"local_id": "GT-065", "original_number": 65, "kind": "P", "entry_kind": "statement_only", "published": false, "review_required": true, "category": "game_theory"}
\label{game:GT-065}
{\sffamily\small\color{muted}
\noindent\textbf{ID: GT-065.}\quad\textbf{Category:} Fair division and resource augmentation.
\quad\textbf{Kind: P.}\quad\textbf{Status:} Open in the cited source for general \(n\).
\par}
\subsection*{Definitions and mathematical statement}
Let \(N=\{1,\ldots,n\}\), \(n\geq 1\), and let \(M\) be a finite set
of goods, with additive nonnegative values \(v_i(S)=\sum_{g\in S}a_{ig}\).
Let \(\Pi_n(M)\) be the set of ordered partitions of \(M\) into \(n\)
parts, including empty parts. Define the original maximin share
\[
 \mu_i=\max_{(P_1,\ldots,P_n)\in\Pi_n(M)}
              \min_{j\in N}v_i(P_j).
\]
Represent copies by bundles \(A_i\subseteq M\) that may overlap.
Each agent may receive at most one instance of each original good.
For \(g\in M\), define its number of additional used copies by
\[
 d_g(A)=\max\bigl\{0,\ |\{i\in N:g\in A_i\}|-1\bigr\}.
\]

\paragraph{Question.}
For every such instance, do there exist \(A_1,\ldots,A_n\subseteq M\)
such that
\[
 \bigcup_{i\in N}A_i=M,\qquad
 \sum_{g\in M}d_g(A)\leq 1,\qquad
 v_i(A_i)\geq\mu_i\quad(i\in N)?
\]
Thus, at most one good is duplicated once, and its two instances must be
assigned to different agents.

\subsection*{Short English statement}
Is one additional copy sufficient to give everyone their original MMS?

\subsection*{Variants and scope boundaries}
The benchmark \(\mu_i\) is computed before augmentation. Recomputing MMS
in the augmented market is a different question. Using at most one copy
handles empty markets and instances needing no duplication. Requiring all
original goods to be used is harmless here: assigning an unused original
good to any agent preserves the displayed inequalities.

\subsection*{Source relationship and status}
This is the first question in [S1, Section 1.2], with the copy convention
of Definition 2.5. Corollary 3.9 answers it positively for three agents.

\subsection*{References}
\begingroup\footnotesize\raggedright
\noindent[S1] Hannaneh Akrami, Siddharth Barman, Alon Eden,
Michal Feldman, Amos Fiat, Yoav Gal-Tzur, Satyanand Rammohan, and
Aditi Sethia,
\emph{Maximin-Share Fairness via Resource Augmentation},
2025; arXiv:2502.09377v3, 2026
(the abstract-page title is \emph{Fair Division via Resource Augmentation}).
\url{https://arxiv.org/html/2502.09377v3}.
Locators: Section 1.2, first and fourth open questions;
Section 2, Definitions 2.1 and 2.5; Corollary 3.9.
\par\endgroup

\subsection*{Common conventions: Source mathematical conventions}

\textbf{Finite games.} For a finite set $A$, $\Delta(A)$ is its probability
simplex. Players choose independent mixed strategies in Nash equilibrium;
correlation is allowed only when explicitly part of the solution concept.
In a game with payoff functions $u_i$ and strategy profile $x$, an additive
$\varepsilon$ Nash equilibrium satisfies
\[
 u_i(x)\geq u_i(x_i',x_{-i})-\varepsilon
 \quad\text{for every player $i$ and every admissible deviation $x_i'$}.
\]
For numerical approximation, payoffs are normalized as locally specified.
An $\varepsilon$ payoff residual is distinct from distance at most
$\varepsilon$ to the set of exact equilibria. Point convergence, convergence
of empirical play, and convergence in distance to an equilibrium set are
also distinct.

\textbf{Computational representations.} Unless stated otherwise, a complexity
question takes rational data with binary encoding; polynomial time is measured
in total bit length. A fixed-accuracy polynomial algorithm is not necessarily
polynomial in $1/\varepsilon$, and neither is necessarily polynomial in
$\log(1/\varepsilon)$. Succinct games and value, demand, or sampling oracles
must declare their interfaces. Exact real solutions require an explicit
representation; an unspecified list of arbitrary real numbers is not a
finite computational output. A claimed algorithmic barrier is meaningful
only for its specified model and reductions.

\textbf{Dynamic games.} Histories, information, observation, and allowable
randomization are part of the model. A stationary strategy depends only on
the current state; a positional strategy in a deterministic graph game
chooses one action at each controlled vertex. Nash equilibrium at an initial
state and equilibrium after every history are different requirements.
An expectation of a pathwise $\liminf$ payoff, a $\liminf$ of expected averages,
a limiting expected average, and a uniform finite-horizon equilibrium are
different criteria. None is silently substituted for another.

\textbf{Allocation and mechanisms.} A complete allocation assigns every item
exactly once unless otherwise stated. Zero-valued goods, negative values,
capacity constraints, feasibility, lotteries, and copies of goods affect
fairness definitions and are specified locally. Dominant-strategy incentive
compatibility, Bayesian incentive compatibility, universal truthfulness,
and truthfulness in expectation impose different inequalities. Whenever
an objective ratio has zero denominator, its convention is stated or those
instances are excluded.

\textbf{Infinite-dimensional models.} Expectations, integrals, stochastic
equations, and optimization objectives require their local regularity,
integrability, filtrations, and admissible domains. A backward--forward
mean-field system is a boundary-value problem; it is not an initial-value
semiflow without an additional construction. Long-time, large-population,
and vanishing-noise limits require an order of limits and a convergence
topology. If a minimum need not be attained, the objective is its infimum
or supremum and the approximation requirement is stated separately.
% END ECONOMICS STATEMENT
\end{document}
